% !TeX root = ../../main.tex
\section{\bilingual{Background and Motivation}{研究背景}}

    \bilingual{
        Branching Random Walk (BRW) is a fundamental model in the theory of stochastic processes that describes the dynamic evolution of populations in space. From a mathematical construction perspective, BRW is not a simple collection of random walks (RW), but rather a complex random system composed of a family of random variables indexed on the Ulam--Harris--Neveu tree vertex space $\mathbb V$.
    }{
        {\kaishu 分支随机游走} (branching random walk, BRW) 是随机过程理论中刻画种群在空间中动态演化的核心模型. ~从数学构造上看,~ BRW 并非{\kaishu 随机游走} (random walk, RW) 的简单集合,~ 而是由一族索引在 {\kaishu Ulam--Harris--Neveu \mycite{Neveu1986} 树节点空间} $\mathbb V$ 上的随机变量构成的复杂随机系统.
    }

    \bilingual{
        Unlike random walks with independent and identically distributed increments, the defining feature of BRW is its significant {\itshape family correlation} or {\itshape hierarchical correlation} among particle displacements. Specifically, for two particles $u,v\in\mathbb V$ at generation $n$, the correlation strength between their spatial positions $V(u)$ and $V(v)$ depends strictly on their geometric distance in the hierarchical tree. If we denote $u\wedge v$ as their {\itshape most recent common ancestor} (MRCA), then the covariance between the two particles is proportional to the length of their shared ancestral path $|u\wedge v|$. This means that particles with closer ``kinship'' have more overlapping trajectories in space, and this nested hierarchical structure endows BRW with deep statistical dependencies.
    }{
        与具有独立同分布位移序列的随机游走不同,~ BRW 的核心特征在于其粒子位移具有显著的{\kaishu 家族相关性} (family correlation) 或{\kaishu 层次相关性} (hierarchical correlation). ~具体而言,~ 对于第 $n$ 代的两个粒子 $u,v\in\mathbb V$,~ 其空间位置 $V(u)$ 与 $V(v)$ 的相关性强度严格取决于它们在层级树中的几何距离. ~若记 $u\wedge v$ 为它们的{\kaishu 最近共祖} (most recent common ancestor, MRCA),~ 则两个粒子的协方差与它们共享的祖先路径长度 $|u\wedge v|$ 成正比. ~这意味着``血缘关系''越近的粒子,~ 其在空间中的运动轨迹重合度越高,~ 这种嵌套的层级结构使得 BRW 表现出极深的统计依赖性.
    }

    \bilingual{
        Due to this hierarchical correlation, BRW has been regarded as a discrete benchmark model for studying {\itshape log-correlated fields}. A log-correlated field is a class of random fields where the covariance between two points grows logarithmically with the inverse of their spatial distance. In a BRW with branching number $d$, the total number of particles at generation $n$ is $N = d^n$, which implies that the depth (hierarchy) of the tree satisfies $n = \log_d N$. Since the maximum variance (i.e., self-correlation) of the displacement between two particles is proportional to the generation $n$, the field intensity (variance) exhibits a logarithmic growth characteristic with respect to the system size $N$ (${\mathrm Var} \sim \log N$). This structure places BRW in the same universality class \mycite{Fyodorov_2012,arguin2016extremalogcorrelatedrandomvariables} as {\itshape two-dimensional Gaussian free fields} \mycite{bramson2015convergencelawmaximumtwodimensional}, {\itshape fluctuations of eigenvalues in random matrices}, and {\itshape behavior of the Riemann $\zeta$ function on the critical line} \mycite{najnudel2018extremevaluesriemannzeta,arguin2022maximarandommodelriemann}.
    }{
        正是由于这种层次相关性,~ BRW 已被视为研究{\kaishu 对数相关场} (log-correlated fields) 的离散基准模型. ~所谓对数相关场,~ 是指一类随机场,~ 其中两点之间的协方差随其空间距离的倒数呈对数级增长. ~在分支数为 $d$ 的 BRW 中,~ 第 $n$ 代的粒子总数 $N = d^n$,~ 这意味着树的深度 (层次) 满足 $n = \log_d N$. ~由于两个粒子位移的最大方差 (即自相关性) 正比于代数 $n$,~ 因此在种群规模为 $N$ 的系统中,~ 场强度 (方差) 相对于系统规模 $N$ 呈现出对数级的增长特征 (${\mathrm Var} \sim \log N$). ~这种结构使得 BRW 与{\kaishu 二维高斯自由场} \mycite{bramson2015convergencelawmaximumtwodimensional}、{\kaishu 随机矩阵特征值涨落}以及{\kaishu 数论中的 Riemann $\zeta$函数在临界线上的行为} \mycite{najnudel2018extremevaluesriemannzeta,arguin2022maximarandommodelriemann} 属于同一个{\kaishu 普适性类} (universality class) \mycite{Fyodorov_2012,arguin2016extremalogcorrelatedrandomvariables}.
    }

    \bilingual{
        Studying the extremal behavior of BRW (such as the position of the extremal particle) is essentially an exploration of how to find a precise probabilistic balance between the ``explosion'' caused by branching growth (exponential sample size) and the ``effective sparsity'' caused by hierarchical correlation. In contrast, in an independent and identically distributed sequence, the extremal behavior is entirely determined by the sample size; however, in BRW, due to the presence of hierarchical correlation, extremal particles often exhibit a {\itshape clustering} characteristic - the emergence of a ``genius'' particle is often accompanied by a collective emergence of its close relatives near the spatial frontier boundary. This correlation significantly alters the order of front advancement, leading to a slower sub-logarithmic delay compared to the independent case \mycite{Bramson1978MaximalDO,hu2009minimalpositioncriticalmartingale,Addario_Berry_2009,A_d_kon_2013,zeitouni2025log}. Therefore, BRW is not only a spatial model but also a window into understanding the universality of extremes in complex correlated systems.
    }{
        研究 BRW 的极值行为 (如极值位置),~ 本质上是在探索如何在分支增长导致的``数量爆炸” (指数级的样本量) 与层次相关性导致的``有效样本稀疏''之间寻找精确的概率平衡. ~作为比较,~ 在独立同分布序列中,~ 极值完全由样本量决定;~ 而在 BRW 中, 由于层次相关性的存在,~ 极值粒子往往呈现出{\kaishu 成簇} (clustering) 的特性——一个``天才''粒子的出现往往伴随着其近亲族群在空间前沿边界附近的集体涌现.~ 这种相关性极大地改变了前沿推进的阶数,~ 导致了比独立情形更慢的次阶对数延迟 \mycite{Bramson1978MaximalDO,hu2009minimalpositioncriticalmartingale,Addario_Berry_2009,A_d_kon_2013,zeitouni2025log}.~ 因此,~ BRW 不仅是一个空间模型,~ 更是理解复杂相关系统中极值普适性的窗口.
    }

    \bilingual{
        As a vibrant branch of stochastic process theory, BRW is not only an ideal mathematical model for studying sequences of random variables with hierarchical correlation, but also a core tool in statistical physics, evolutionary biology, and population genetics for characterizing spatiotemporal evolutionary dynamics. In non-equilibrium statistical physics, BRW has a profound duality with the nonlinear partial differential equation - the Fisher--KPP (Kolmogorov--Petrovsky--Piskunov) equation. The classical work of McKean~\mycite{mckean1975application} proved that the distribution function of the continuous version of BBM (branching Brownian motion) corresponds exactly to the solution of the KPP equation, thus linking the front propagation speed of the population to the intrinsic speed $v_\infty$ of the traveling wave solution. However, the deterministic KPP equation describes diffusion in a continuous medium, neglecting the stochastic fluctuations arising from the discrete {\itshape atomicity} of real physical systems. To simulate this microscopic effect, the $N$-branching random walk ($N$-BRW) model was introduced. Brunet and Derrida~\mycite{PhysRevE.56.2597} heuristically pointed out that under the constraint of finite population size $N$, there is a tiny ``cutoff effect'' on the front position; due to the exponential amplification effect of branching processes on front perturbations, this microscopic-level change ultimately leads to a significant $(\log N)^{-2}$ order correction in macroscopic wave speed.
    }{
        作为随机过程理论中极具生命力的分支,~ BRW 不仅是数学上研究具有层次相关性随机变量序列的理想模型,~ 更是统计物理、演化生物学及群体遗传学中刻画时空演化动力学的核心工具. ~在非平衡态统计物理中,~ BRW 与非线性偏微分方程——Fisher--KPP (Kolmogorov--Petrovsky--Piskunov) 方程有着深刻的对偶性. ~McKean~\mycite{mckean1975application} 的经典工作证明了BBW的连续版本{\kaishu 分支Brownian运动} (branching Brownian motion, BBM) 的分布函数恰好对应 KPP 方程的解,~ 从而将种群的前沿传播速度与行波解的本征速度 $v_\infty$ 联系起来.~ 然而,~ 确定性的 KPP 方程描述的是连续介质的扩散,~ 忽略了现实物理系统由于离散的{\kaishu 原子性} (atomicity) 而产生的随机涨落. ~为了模拟这种微观效应, {\kaishu $N$-分支随机游走} ($N$-branching random walk, $N$-BRW) 模型应运而生. ~Brunet 与 Derrida~\mycite{PhysRevE.56.2597} 启发式地指出,~ 在有限种群规模 $N$ 的限制下,~ 前沿位置会产生一个微小的``截断效应'';~ 由于分支过程对前沿扰动具有指数级的放大作用,~ 这一微观层面的改变最终导致宏观波速产生显著的 $(\log N)^{-2}$ 量级的修正.
    }

    \begin{figure}[H]
        \centering
        \begin{tikzpicture}[
            xscale=1.8, % 修改这里：控制横向（代数 n）的间距
            yscale=0.8, % 修改这里：控制纵向（空间位置 V）的伸展程度
            % 存活节点样式：蓝色实心
            survivor/.style={circle, fill=blue!70!black, inner sep=1.5pt},
            % 被灭杀节点样式：灰色空心
            killed/.style={circle, draw=black, fill=white, inner sep=1.5pt},
            % 连线样式
            live-edge/.style={thick, blue!70!black},
            killed-edge/.style={dashed, black},
            % 叉叉样式
            cross/.style={cross out, draw=red!70!black, inner sep=0pt, minimum size=4pt, thick}
        ]

            % --- 坐标轴提示 ---
            \draw[->, >=Stealth, black, thick] (0, -3) -- (0, 2.5);
            \draw[-, black, thick, dashed] (2, -3) -- (2, 2.5);
            \draw[-, black, thick, dashed] (4, -3) -- (4, 2.5);

            % --- 第 0 代 (n=0) ---
            \node[survivor] (root) at (0, 0) {};
            \node[anchor=south] at (0, 2.7) {$n=0$};

            % --- 第 1 代 (n=1, 产生4个分叉) ---
            % 设定 4 个子代的位置，y轴越小越靠下
            \node[survivor] (g1-1) at (2, -1.6) {}; % 存活 1
            \node[survivor] (g1-2) at (2, -0.3) {}; % 存活 2
            \node[survivor] (g1-3) at (2, 0.5) {};  % 存活 3
            \node[killed]   (g1-4) at (2, 1.8) {};  % 被淘汰 (第4个)

            % 第 1 代连线
            \draw[live-edge] (root) -- (g1-1);
            \draw[live-edge] (root) -- (g1-2);
            \draw[live-edge] (root) -- (g1-3);
            \draw[killed-edge] (root) -- (g1-4);
            % \node[cross] at (g1-4) {};

            \node[anchor=south] at (2, 2.7) {$n=1$};

            % --- 第 2 代 (n=2, 存活的3个父节点各产生4个子代，共12个) ---

            % 为了图表清晰，我们假设第2代的 3 个存活粒子主要集中在最下方的分支

            % 1. 来自 g1-1 的子代 (假设位置很靠下，存活较多)
            \node[survivor] (g2-1-1) at (4, -2.6) {};
            \node[survivor] (g2-1-2) at (4, -1.5) {};
            \node[killed]   (g2-1-3) at (4, 0.4) {};
            \node[killed]   (g2-1-4) at (4, 0.9) {};
            \foreach \i in {1,2} \draw[live-edge] (g1-1) -- (g2-1-\i);
            \foreach \i in {3,4} {\draw[killed-edge] (g1-1) -- (g2-1-\i); }% \node[cross] at (g2-1-\i) {};}

            % 2. 来自 g1-2 的子代 (只有1个进入全局前3)
            \node[survivor] (g2-2-1) at (4, -1.8) {};
            \node[killed]   (g2-2-2) at (4, -0.4) {};
            \node[killed]   (g2-2-3) at (4, -0.8) {};
            \node[killed]   (g2-2-4) at (4, 0.4) {};
            \draw[live-edge] (g1-2) -- (g2-2-1);
            \foreach \i in {2,3,4} {\draw[killed-edge] (g1-2) -- (g2-2-\i); }%\node[cross] at (g2-2-\i) {};}

            % 3. 来自 g1-3 的子代 (全部被淘汰，因为全局排名在3之后)
            \foreach \y/\i in {-0.6/1, -0.3/2, 1.2/3, 1.9/4} {
                \node[killed] (g2-3-\i) at (4, \y) {};
                \draw[killed-edge] (g1-3) -- (g2-3-\i);
                % \node[cross] at (g2-3-\i) {};
            }

            \node[anchor=south] at (4, 2.7) {$n=2$};

        \end{tikzpicture}
        \caption{\bilingual{Illustration of a \texorpdfstring{$4$}{4}-branching \texorpdfstring{$3$}{3}-branching random walk}{\texorpdfstring{$4$}{4}-分叉的\texorpdfstring{$3$}{3}-分支随机游走示例图}}
    \end{figure}

    \bilingual{
        More precisely,~ let $\Xi$ be a point process on $\mathbb R$,~ referred to as the {\itshape branching law}. Starting with $\lfloor N^a\rfloor$ particles at position $0$ at time $0$ (where $a\in[0,1]$),~ at each generation:\, each particle independently produces offspring according to an independent copy of $\Xi$, with displacements given by the atoms of that copy; then, among all offspring, only the $N$ leftmost particles (i.e., those with the smallest positions) are retained for the next generation. This process is called the {\itshape $N$-branching random walk ($N$-BRW)}, denoted by $\mathcal V^N$ or $\mathcal V^{N|a}$. The minimum position at generation $n$ is denoted by $\mathcal M_n^N:=\min\mathcal V^N_n$.
    }{
        更确切地说,~ 记 $\Xi$ 为 $\mathbb R$ 上的一个点过程,~ 称为{\kaishu 分支律}. 在初始时刻 $0$,~ 有 $\lfloor N^a\rfloor$ ($a\in[0,1]$) 个粒子从位置 $0$ 出发;~ 在每一代中:\, 每个粒子根据 $\Xi$ 的一个独立副本生成后代,~ 子代位移由该副本的各个原子给出;~ 然后,~ 在所有子代中,~ 仅保留位置最靠前的 $N$ 个粒子进入下一代. 这一过程称为{\kaishu $N$-分支随机游走 ($N$-BRW)}, 记作 $\mathcal V^N$ 或 $\mathcal V^{N|a}$; 第 $n$ 代的最左侧位置记作 $\mathcal M_n^N:=\min\mathcal V^N_n$.
    }

    \bilingual{
        The significance of Brunet--Derrida's discovery lies in its challenge to the conventional wisdom in statistical physics that ``noise can be neglected in the large $N$ limit''. It revealed a non-trivial {\itshape scaling law}: even with an infinitesimal microscopic cutoff, due to the logarithmic sensitivity of front dynamics to discreteness, a significant macroscopic velocity shift can still occur. This finding forced the physics community to re-examine the limitations of deterministic differential equations in describing discrete living systems and physical fronts, and established the new universal paradigm of {\itshape noise-induced front shift}.
    }{
        Brunet--Derrida 的发现之所以具有里程碑意义,~ 是因为它打破了传统统计物理中``大$N$极限下噪声可忽略''的固有思维. 它揭示了一种非平凡的{\kaishu 标度律} (scaling law):~ 即使微观截断极小,~ 由于前沿动力学对离散性的对数级敏感性,~ 宏观速度仍会产生显著的偏移. ~这一发现迫使物理学界重新审视确定性微分方程在描述离散生命系统与物理前沿时的局限性,~ 并确立了{\kaishu 噪声诱导前沿偏移} (noise-induced front shift) 这一全新的普适性范式.
    }

    \bilingual{
        This discovery not only reshaped the understanding of nonlinear wave dynamics in statistical physics, but also exhibited profound universality significance in evolutionary biology and the physics of disordered systems. In evolutionary biology, $N$-BRW has become a standard model for simulating natural selection under limited resources. In this context, the spatial position $\mathcal V(u)$ of a particle is endowed with the biological meaning of {\itshape individual fitness}~\mycite{tsimring1996rna}. Although populations tend to expand in supercritical environments, the finite {\itshape environmental carrying capacity} $N$ enforces a selection pressure of ``survival of the fittest'': in each generation of reproduction, only the $N$ individuals with the highest fitness can survive and reproduce. This truncation mechanism not only fundamentally determines the macroscopic rate of adaptive evolution but also profoundly reshapes the {\itshape genealogical structure} of populations~\mycite{Brunet_2007,Berestycki_2013}. Through in-depth study of $N$-BRW, mathematical biologists have been able to quantitatively explain how populations maintain genetic diversity through front diffusion under natural selection pressure and reveal the dynamic characteristics of advantageous mutations in population expansion~\mycite{desai2007beneficial,PMID:23269838}. This interdisciplinary relevance indicates that $N$-BRW and its related models have become an important bridge connecting microscopic randomness, nonlinear dynamics, and macroscopic properties of complex systems.
    }{
        这一发现不仅重塑了统计物理关于非线性波动的认知,~ 更在演化生物学与无序系统物理学中展现出深远的普适性意义. ~在演化生物学领域,~ $N$-BRW 已成为模拟受限资源下自然选择的标准模型.~ 在该语境下,~ 粒子的空间位置 $\mathcal V(u)$ 被赋予了{\kaishu 个体适应度} (fitness)~\mycite{tsimring1996rna} 的生物学含义. ~虽然在超临界环境下种群具有扩张趋势,~ 但有限的{\kaishu 环境承载力} (carrying capacity) $N$ 强制执行了``优胜劣汰''的选择压力:~ 在每一代繁殖中,~ 仅有适应度最高的 $N$ 个个体才能得以生存并繁衍. ~这一截断机制不仅从根本上决定了适应性进化的宏观速率,~ 更深刻地重塑了种群的{\kaishu 系谱结构} (genealogical structure)~\mycite{Brunet_2007,Berestycki_2013}. ~通过对 $N$-BRW 的深入研究, ~生物数学家得以定量解释在自然选择压力下, ~种群如何通过前沿扩散维持遗传多样性,~ 并揭示优势突变在种群扩张中的动力学特征~\mycite{desai2007beneficial,PMID:23269838}. ~这种跨学科的关联性表明, ~$N$-BRW 及其相关模型已成为连接微观随机性、非线性动力学与复杂系统宏观性质的重要桥梁.
    }
